
\Needspace{12\baselineskip}
\subsubsection*{Naučnoistraživački rad -- \ClanJedan}

\textbf{Naslov rada:}
\emph{A supervised machine learning algorithm for detecting
and predicting fraud in credit card transactions}.\\

\textbf{Autori:}
Jonathan Kwaku Afriyie, Kassim Tawiah, Wilhemina Adoma Pels,
Sandra Addai-Henne, Harriet Achiaa Dwamena, Emmanuel Odame Owiredu,
Samuel Amening Ayeh i John Eshun.\\

\textbf{Izvor:}
\emph{Decision Analytics Journal}, vol. 6, članak 100163,
2023. \cite{afriyie2023fraud}

\begin{table}[H]
\centering
\caption{Sažeti pregled naučnoistraživačkog rada Afriyie i saradnika (2023)}
\label{tab:afriyie-pregled}
\small
\renewcommand{\arraystretch}{1.22}
\begin{tabularx}{\linewidth}{@{}p{0.25\linewidth}Y@{}}
\toprule
\textbf{Stavka} & \textbf{Opis} \\
\midrule

Problem &
Detekcija prevarnih transakcija kreditnim karticama kao problem
binarne klasifikacije, pri čemu se razlikuju legitimne i
prevarne transakcije. \\

Skup podataka &
\emph{Credit Card Transactions Fraud Detection Dataset}
(Kaggle), simuliran pomoću alata Sparkov.
U radu se navodi 555\,719 transakcija i 23 kolone
u početnom skupu podataka. \\

Korištene metode &
Logistička regresija (Logistic Regression),
stablo odlučivanja (Decision Tree) i slučajna šuma
(Random Forest). Priprema podataka uključuje
skaliranje i nasumično smanjivanje većinske klase
(undersampling). \\

Metrike &
Tačnost (accuracy), preciznost (precision),
osjetljivost (recall/sensitivity), specifičnost,
F1-mjera i ROC-AUC. \\

Najbolji rezultat &
\textbf{Random Forest}: tačnost \textbf{96\%},
ROC-AUC \textbf{98,9\%}.
Prijavljene vrijednosti: recall 0,97,
precision 0,09 i F1-mjera 0,17. \\

\bottomrule
\end{tabularx}
\end{table}

\textbf{Cilj i istraživački problem.}

Autori istražuju koliko uspješno različiti algoritmi nadziranog
mašinskog učenja razlikuju prevarne od legitimnih transakcija.
Osnovni cilj jeste uporediti performanse tri klasifikacijska
algoritma i utvrditi koji je najpogodniji za analizirani
skup podataka.

Dodatno istražuju obrasce povezane s iznosom i vremenom
transakcije, kategorijom trgovca te pojedinim
karakteristikama korisnika.

\textbf{Skup podataka i njegova nebalansiranost.}

Rad koristi javno dostupan Kaggle skup simuliranih transakcija,
obavljenih tokom 2020. godine u Sjedinjenim Američkim Državama.

Skup obuhvata numeričke i kategoričke atribute,
a ciljna varijabla predstavlja status prevare.

Nakon postupaka obrade, u Tabeli 4 rada prikazano je
482\,672 legitimnih i 2\,135 prevarnih transakcija,
odnosno približno 99,6\% i 0,4\%.

Takva raspodjela predstavlja izraženu nebalansiranost klasa,
pa sama ukupna tačnost nije dovoljna za procjenu
uspješnosti klasifikatora.

\textbf{Priprema podataka i korištene metode.}

Autori navode čišćenje podataka, transformaciju numeričkih
atributa na raspon od 0 do 1 i primjenu metode
\emph{random undersampling}.

Ova metoda smanjuje broj primjera većinske klase kako bi
se ublažila pristrasnost modela prema legitimnim transakcijama.

Na Slici 1 rada prikazan je uravnotežen uzorak s
1\,708 prevarnih i 1\,707 legitimnih transakcija.

Za klasifikaciju se zatim koriste logistička regresija,
jedno stablo odlučivanja i Random Forest,
koji kombinuje veći broj stabala.

\textbf{Poređenje eksperimentalnih rezultata.}

U Tabeli~\ref{tab:afriyie-rezultati} prikazane su vrijednosti
koje su autori objavili u Tabeli 11 i uz ROC krive (Slika 16).

\begin{table}[H]
\centering
\caption{Poređenje rezultata klasifikacijskih algoritama iz rada}
\label{tab:afriyie-rezultati}
\small
\setlength{\tabcolsep}{4pt}
\renewcommand{\arraystretch}{1.18}

\begin{tabularx}{\linewidth}{@{}Y c c c c c c@{}}
\toprule

\textbf{Model} &
\textbf{Acc.} &
\textbf{Prec.} &
\textbf{Recall} &
\textbf{F1} &
\textbf{Spec.} &
\textbf{AUC} \\

\midrule

Decision Tree &
0,92 & 0,05 & 0,93 & 0,09 & 0,92 & 0,945 \\

Logistic Regression &
0,92 & 0,04 & 0,76 & 0,08 & 0,92 & 0,872 \\

\textbf{Random Forest} &
\textbf{0,96} &
\textbf{0,09} &
\textbf{0,97} &
\textbf{0,17} &
\textbf{0,96} &
\textbf{0,989} \\

\bottomrule
\end{tabularx}
\end{table}

\textbf{Tumačenje rezultata.}

Random Forest je postigao najveće vrijednosti svih
prikazanih metrika.

Njegov ROC-AUC od 0,989 ukazuje na veoma dobru sposobnost
razlikovanja klasa pri promjeni klasifikacijskog praga.

Međutim, precision od samo 0,09 znači da je tek približno
9\% transakcija označenih kao prevarne zaista bilo prevarno.

Iz objavljene konfuzijske matrice (Tabela 7 rada)
vidi se da je model prepoznao 409 prevara i propustio 18,
ali je 4\,052 legitimne transakcije pogrešno označio
kao prevarne.

Zato visoka tačnost i ROC-AUC ne znače da model proizvodi
malo lažnih alarma.

Niska F1-mjera (0,17) dodatno ukazuje na neravnotežu
između preciznosti i osjetljivosti modela.

\textbf{Dodatna zapažanja iz istraživačke analize.}

Autori izvještavaju da se veći broj prevarnih transakcija
u njihovom skupu javlja između 22:00 i 04:00 GMT te
posebno razmatraju starosnu grupu korisnika iznad
60 godina i kategoriju kupovine.

Budući da su podaci simulirani, ova zapažanja
predstavljaju obrasce uočene u posmatranom skupu,
a ne dokazane zakonitosti stvarnog finansijskog poslovanja.

\textbf{Zaključak i značaj za istraživanje.}

Istraživanje upućuje na to da Random Forest može biti
uspješan izbor za detekciju prevarnih transakcija,
posebno prema ROC-AUC i osjetljivosti.

Istovremeno pokazuje da izrazita nebalansiranost klasa
zahtijeva pažljivu evaluaciju i praćenje lažnih alarma.

Nalazi rada predstavljaju osnovu za očekivanje da će
metode zasnovane na stablima ostvarivati konkurentne rezultate,
ali i da će precision, recall i F1-mjera biti važnije
za praktičnu interpretaciju nego sama ukupna tačnost.
